\documentclass[../main.tex]{subfiles}

\begin{document}

\subsection{Applying the Roofline Model: Block Tiling and Thread Coarsening}

\begin{frame}
  \frametitle{Driving Question}

  \begin{alertblock}{How can we incrase AI?}
    The roofline model tells us that increasing AI should increase throughput.
    \textbf{But how can we actually increase AI?}
  \end{alertblock}

  \begin{enumerate}
    \item \pause Utilizing SMEM more aggressively.
    \item \pause Loading larger blocks into SMEM.
    \item \pause Reusing values in SMEM to compute multiple values of \(C\).
  \end{enumerate}
\end{frame}

\begin{frame}
  \frametitle{Generalized Tiling: Loading larger blocks.}
  Our tiled kernel \textbf{hardcoded} a square \(32 \times 32\) block size. Each
  thread in the block load one value into SMEM. We can parameterize the tiling
  with \(BM\), \(BN\), \(BK\).
\end{frame}

\begin{frame}[fragile]
  \frametitle{Basic tiling: one thread, one element from each tile}

  \begin{columns}[T,onlytextwidth]
    \begin{column}{0.61\textwidth}
      \begin{lstlisting}[
        basicstyle=\ttfamily\fontsize{7.2}{8.0}\selectfont,
        breaklines=false,
        escapeinside={(*@}{@*)}
      ]
__shared__ float As[BK][BK];
__shared__ float Bs[BK][BK];

for (int k = 0; k < K; k += BK) {
  const int Arow = y;
  const int Acol = k + xt;
(*@\tikz[remember picture,overlay]\coordinate (basic-load-start);@*)  if (Arow < M && Acol < K)
    As[yt][xt] = A[Arow * K + Acol];
  else
    As[yt][xt] = 0.0;

  const int Brow = k + yt;
  const int Bcol = x;
  if (Brow < K && Bcol < N)
    Bs[yt][xt] = B[Brow * N + Bcol];(*@\tikz[remember picture,overlay]\coordinate (basic-load-right);@*)
  else
    Bs[yt][xt] = 0.0;(*@\tikz[remember picture,overlay]\coordinate (basic-load-end);@*)

  __syncthreads();
  // compute with As and Bs ...
}
      \end{lstlisting}
    \end{column}

    \begin{column}{0.36\textwidth}
      \begin{block}{Direct mapping}
        \small
        The thread coordinates directly select the shared-memory cells:
        \[
          (y_t,x_t) \longrightarrow
          \begin{cases}
            A_s[y_t][x_t]\\
            B_s[y_t][x_t]
          \end{cases}
        \]
      \end{block}
      \begin{alertblock}{Limitation}
        \small
        This only works because the block and both shared-memory tiles are all
        $BK\times BK$.
      \end{alertblock}
    \end{column}
  \end{columns}

  \begin{tikzpicture}[remember picture,overlay]
    \draw[orange!85!black,very thick,rounded corners]
    ($(basic-load-start)+(-1mm,1.2ex)$) rectangle
    ($(basic-load-right |- basic-load-end)+(1mm,-0.5ex)$);
  \end{tikzpicture}
\end{frame}

\begin{frame}[fragile]
  \frametitle{General tiling: threads cooperatively load each tile}

  \begin{columns}[T,onlytextwidth]
    \begin{column}{0.64\textwidth}
      \begin{lstlisting}[
        basicstyle=\ttfamily\fontsize{6.6}{7.3}\selectfont,
        breaklines=false,
        escapeinside={(*@}{@*)}
      ]
const int tid = yt * BN + xt;
const int numThreads = BM * BN;

__shared__ float As[BM * BK];
__shared__ float Bs[BK * BN];
for (int k = 0; k < K; k += BK) {
(*@\tikz[remember picture,overlay]\coordinate (general-load-start);@*)  for (int id = tid; id < BM * BK;
       id += numThreads) {
    const int yid = id / BK;
    const int xid = id % BK;
    const int Arow = yb + yid;
    const int Acol = k + xid;
    As[id] = (Arow < M && Acol < K)
           ? A[Arow * K + Acol] : 0.0;
  }
  for (int id = tid; id < BK * BN;
       id += numThreads) {
    const int yid = id / BN;
    const int xid = id % BN;
    const int Brow = k + yid;
    const int Bcol = xb + xid;
    Bs[id] = (Brow < K && Bcol < N)
           ? B[Brow * N + Bcol] : 0.0;(*@\tikz[remember picture,overlay]\coordinate (general-load-right);@*)
  }(*@\tikz[remember picture,overlay]\coordinate (general-load-end);@*)
  __syncthreads();
  // Compute with As and Bs ...
}
      \end{lstlisting}
    \end{column}

    \begin{column}{0.33\textwidth}
      \begin{block}{Thread-strided loading}
        \small
        Flatten each tile. Thread $tid$ loads
        \[
          id=tid+jT
        \]
        for $j=0,1,\ldots$, where $T=BM\cdot BN$ is the thread count.
      \end{block}
      \begin{alertblock}{Key difference}
        \small
        A thread may load zero, one, or several values. Loading no longer
        assumes $BM=BN=BK$.
      \end{alertblock}
    \end{column}
  \end{columns}

  \begin{tikzpicture}[remember picture,overlay]
    \draw[orange!85!black,very thick,rounded corners]
    ($(general-load-start)+(-1mm,1.2ex)$) rectangle
    ($(general-load-right |- general-load-end)+(1mm,-0.5ex)$);
  \end{tikzpicture}
\end{frame}

\begin{frame}
  \frametitle{Visualization}
  \begin{overlayarea}{\linewidth}{0.58\textheight}
    \centering
    \only<1>{
      \includegraphics[width=0.9\linewidth,height=0.58\textheight,keepaspectratio]{tiled_general/basic_thread_assignment.png}
    }
    \only<2>{
      \includegraphics[width=0.9\linewidth,height=0.58\textheight,keepaspectratio]{tiled_general/general_thread_assignment.png}
    }
  \end{overlayarea}
\end{frame}

\begin{frame}
  \frametitle{Thread coarsening: Increasing AI}
  \begin{itemize}
    \item Basic tiling assigns $1$ element of the output to each thread.
    \item We can \emph{coarsen} each thread, and assign a $TM \times TN$ portion of the output to each thread.
    \item We load larger chunks into SMEM, increasing data reuse and AI.
    \item One loaded $A$ value can update several columns; one loaded $B$ value can update several rows.
  \end{itemize}

  \centering
  \includegraphics[width=0.8\linewidth]{data_use_thread_coarsening.png}
\end{frame}

\begin{frame}
  \frametitle{Coarsened SGEMM visualization}
  \centering
  \only<1>{
    \includegraphics[height=.7\textheight]{tiled_steps_coarsened/step_0.png}
  }
  \only<2>{
    \includegraphics[height=.7\textheight]{tiled_steps_coarsened/step_1.png}
  }
\end{frame}

% Unnecessary
% \begin{frame}
%   \frametitle{Coarsening factors are a tradeoff}
%   \begin{columns}[T,onlytextwidth]
%     \begin{column}{0.52\textwidth}
%       \begin{block}{Two levels of tiling}
%         \begin{itemize}
%           \item A block computes a $BM\times BN$ output tile while advancing
%             through $K$ in steps of $BK$.
%           \item Each thread computes a $TM\times TN$ micro-tile in registers.
%           \item Thread count:
%             $\frac{BM}{TM}\frac{BN}{TN}$.
%         \end{itemize}
%       \end{block}
%     \end{column}
%     \begin{column}{0.45\textwidth}
%       \begin{alertblock}{Why not keep increasing them?}
%         Larger tiles improve reuse and arithmetic intensity, but consume more
%         registers and shared memory. Eventually occupancy falls and extra
%         parallelism or reuse no longer helps.
%       \end{alertblock}
%     \end{column}
%   \end{columns}

% \vspace{1mm}
% \centering
% \small
% \begin{tabular}{c|c|c|c|c}
%   Config. & $BM\!\times\!BN\!\times\!BK$ & $TM\!\times\!TN$ & AI & Throughput\\ \hline
%   Basic tiled & $64\times64\times8$ & $4\times4$ & 15.9 & 13.75 TFLOP/s\\
%   C4 & $128\times128\times16$ & $8\times8$ & 31.5 & \textbf{19.54 TFLOP/s}\\
%   C5 & $128\times256\times8$ & $8\times8$ & 41.8 & 17.79 TFLOP/s
% \end{tabular}
% \end{frame}

% Stop using latency as a metric
% \begin{frame}
%   \frametitle{Basic tiled vs. block-tiled latency}
%   \centering
%   \includegraphics[width=0.9\textwidth]{generated/basic_block_tiled_latency.png}
% \end{frame}

\begin{frame}
  \frametitle{Higher-AI tiles tend to raise throughput}
  \centering
  \includegraphics[width=0.9\textwidth]{generated/block_tiled_throughput.png}
\end{frame}

\begin{frame}
  \frametitle{Tiling configurations on the roofline}
  \centering
  \includegraphics[width=0.8\textwidth]{generated/blocktiling_config_roofline.png}

  Reminder that the Roofline Model only tells us the upper bound!
\end{frame}

\begin{frame}
  \frametitle{Higher AI does not guarantee higher throughput}
  \centering
  \includegraphics[width=0.88\textwidth]{generated/higher_intensity_config_roofline.png}

  The roofline is an upper bound, not a performance prediction.
\end{frame}

\begin{frame}
  \frametitle{We have only scratched the surface: cuBLAS}
  \begin{itemize}
    \item cuBLAS is NVIDIA's architecture-tuned linear algebra library and the
      practical performance target for our SGEMM kernels.
    \item At $n=8192$, cuBLAS reaches $26.96$ TFLOP/s, while the
      highest-AI selected configuration (C3) reaches $13.71$ TFLOP/s, or
      $50.8\%$ of cuBLAS's throughput.
    \item cuBLAS completes the multiplication in $40.8$ ms versus $80.2$ ms
      for C3, about $1.97\times$ faster.
  \end{itemize}
  \begin{block}{The remaining gap}
    Production kernels combine deeper register and warp tiling, vectorized
    memory access, pipelined loads, and architecture-specific instruction
    scheduling.
  \end{block}
\end{frame}

% Stop using latency
% \begin{frame}
%   \frametitle{Basic tiled vs. block-tiled latency vs. cuBLAS}
%   \centering
%   \includegraphics[width=0.9\textwidth]{generated/cublas_block_tiled_latency.png}
% \end{frame}

\begin{frame}
  \frametitle{cuBLAS remains the performance target}
  \centering
  \includegraphics[width=0.9\textwidth]{generated/cublas_block_tiled_throughput.png}
\end{frame}

\begin{frame}
  \frametitle{cuBLAS on the roofline}
  \centering
  \includegraphics[width=0.94\textwidth]{generated/cublas_blocktiling_config_roofline.png}

  {\tiny cuBLAS arithmetic intensity is estimated from the profiled
  $256\times128$ output-tile shape.}
\end{frame}

\begin{frame}
  \frametitle{We've come a long way}
  \centering
  \includegraphics[width=0.9\textwidth]{generated/all_implementations_latency.png}
\end{frame}

% Unnecessary
% \begin{frame}
%   \frametitle{All tested tiling configurations}
%   \centering
%   \includegraphics[width=0.94\textwidth]{generated/all_blocktiling_config_roofline.png}
% \end{frame}

\end{document}
